\documentclass[12pt]{article}
\usepackage[utf8]{inputenc}
\usepackage[T1]{fontenc}
\usepackage{hyperref}
\title{Savosavo Language Resources for K-12 Curriculum Development}
\author{Compiled by Claude Sonnet 5.5 (claude-sonnet-5-5) for Andr\'as Kornai}
\date{}
\begin{document}
\maketitle

\section{Introduction}
Savosavo (ISO 639-3 svs; Glottocode savo1255) is a non-Austronesian language of Savo Island, Solomon Islands, with about 2,400 speakers. It is in the Central Solomons family. The UNESCO Atlas lists it as Definitely Endangered. It is not known to be taught in schools. The documentation is substantial, mainly from the DoBeS project and Claudia Wegener's work.

\section{Primary Educational Resources}
The DoBeS project is reported to have produced community materials, including a storybook and a literacy primer. I did not find a downloadable copy of either.
\begin{itemize}
\item DoBeS Savosavo project page: \url{https://dobes.mpi.nl/?p=1259}
\end{itemize}

\section{Linguistic and Reference Materials}
\begin{itemize}
\item Claudia Wegener, \emph{A Grammar of Savosavo} (De Gruyter Mouton, 2012). It includes glossed texts and a lexeme list: \url{https://www.degruyterbrill.com/document/doi/10.1515/9783110289657/html}
\item Savosavo language (Wikipedia): \url{https://wikipedia.com/wiki/ISO_639:svs}
\end{itemize}
I did not find a stand-alone Savosavo dictionary.

\section{Religious and Community Resources}
\begin{itemize}
\item Global Recordings Network, Savosavo audio of an evangelistic kind: \url{https://globalrecordings.net/en/language/svs}
\end{itemize}

\section{Digital and Audio Resources}
\begin{itemize}
\item DoBeS poster on the Savosavo documentation: \url{https://dobes.mpi.nl/posters/print/oceania-savosavo-print.pdf}
\item DoReCo corpus entry for Savosavo (CC BY): \url{https://doreco.huma-num.fr/languages/savo1255}
\item TEITOK search interface for the DoReCo Savosavo data: \url{https://teitok.huma-num.fr/doreco/savo1255/index.php}
\item PARADISEC, Darrell Tryon collection item DT1-163: \url{https://catalog.paradisec.org.au/collections/DT1/items/163}
\item PARADISEC, Capell recordings item AC1-013, which includes Savosavo among several languages: \url{https://catalog.paradisec.org.au/repository/AC1/013}
\end{itemize}

\section{Curriculum Development Context}
Savosavo is spoken by all in the community, which makes it a good candidate for mother-tongue teaching. The grammar, the DoReCo and DoBeS recordings and the primer and storybook produced by the project could serve as a base. There is no known school programme, and some archive access requires registration.

\end{document}